\chapter*{Abstract}
\addcontentsline{toc}{chapter}{Abstract}
Requests, commitments and deadlines in workplace email and chat are rarely stated as explicit to-do items. They appear as polite questions (\emph{``Could you send me the figures?''}), as promises (\emph{``We'll cover this tomorrow''}), as relative dates (\emph{``by Friday''}) and as references to earlier work (\emph{``once the dataset is cleaned''}). Text2Task is a system that reads such messages and turns them into structured, reviewable task proposals: who should do what, by when, and after which other task.

The project combines two fine-tuned DistilBERT models with a substantial deterministic layer. The entity model labels TASK, PERSON and TEMPORAL spans with BIO tags; the relation model classifies task--person and task--date pairs into ASSIGNED\_TO, DUE\_ON and SCHEDULED\_ON using typed entity markers. Around them sit a rule-based refinement pipeline (span repair, temporal detection, grammar-based action-phrase recovery and clause-level dependency extraction), a temporal normaliser, owner resolution, an uncertainty and flagging layer, a cue-based reconciler that updates existing tasks instead of duplicating them, and a dependency graph that refuses cycles. A local web application (FastAPI and React) presents every output as a suggestion that only the person it affects may accept, and keeps an append-only history.

All models were trained on data derived from the MailEx email event-extraction corpus. Because no human-verified gold labels exist, every number in this report measures agreement with provisional SILVER or MailEx-derived labels, or the outcome of developer-written functional checks. On a deduplicated development set of 349 emails, the best entity model reaches a lenient (IoU $\ge 0.5$) F1 of 0.555 and a strict F1 of 0.390, against 0.401 and 0.260 for a rule baseline. The checkpoint used by the application, trained on a smaller SILVER set, is much weaker on TASK spans (strict F1 0.004). The relation model reaches a micro F1 of 0.764 (95\% CI 0.70--0.84) on reference spans, against 0.527 for a rule linker. On 30 held-out chat scenarios, the hybrid configuration passes 23 cases strictly, against 7 for the raw neural models and 22 for rules alone; task recall rises from 0.09 to 0.91.

The central finding is that, at this scale and label quality, the deterministic layer does most of the work on short chat messages, while the neural models contribute most on longer email text, especially for names and dates. The report documents the full development history, including unsuccessful iterations, a train/development leakage problem found during evaluation, and the limitations that prevent any claim of final accuracy.
